\subsection{计算分数选择的严格 CPU 参考}
\label{sec:7-strict-selection}

全部 55,296 个分组查询实例的主码顺序与 CPU 参考精确一致。单 block 始终满足距离顺序，56 个集合差异均限于截断并列；全部 69,120 个实例对已保留精排分数的最终选择均精确（\cref{tab:7-strict-selection}）。

严格 CPU 检查对 GPU 生成的分数按 $(\primarydist,\mathrm{position})$ 排序。导出器遵循扫描布局和累加顺序，有效返回候选的导出分数逐位匹配。源码比较发现相对表扫描版本只增加导出。270 个配置覆盖 Vogue、arXiv 和 OpenAI-3072，probes 为 1、8、128，$\alpha\in\{4,100,200\}$、$\gscan\in\{0,1,2,4,8\}$ 以及两种阈值设置。

\suppressfloats[t]
\begin{table}[!htbp]
\centering
\caption[严格计算分数选择检查]{严格选择检查 CPU 参考检查。集合与顺序相等使用完整主码键。有效计数检查 $\min(c,n_c)$ 个保留候选；最终相等按 $(\refineddist,\mathrm{ID})$ 对保留的计算精排分数排序，包含其 FP32 bits。计数为查询实例，同一查询跨配置复用。}
\label{tab:7-strict-selection}
\footnotesize\setlength{\tabcolsep}{3pt}
\begin{tabular}{@{}lrrrrr@{}}
\toprule
路径 & 查询数 & 集合／顺序精确 & 仅并列 & 有效计数 & 最终精确 \\
\midrule
$\gscan\geq1$ & 55,296 & 55,296 / 55,296 & 0 & 55,296 & 55,296 \\
$\gscan=0$ & 13,824 & 13,768 / 8,758 & 56 & 13,824 & 13,824 \\
\bottomrule
\end{tabular}

\end{table}

全部 19,120 个候选不足案例的有效候选数正确。这些结果检验了 \cref{lem:5-buffer,prop:5-local-coverage,prop:5-shared-bound} 在记录案例中的表现。最终选择相等针对已保留的计算分数，而非原始空间精确近邻；sanitizer 覆盖限于其记录子集（\cref{sec:a-t4-t5}）。
